\section{Conclusion}
\label{sec:conclusion}

We presented a compute-matched study of test-time search for hierarchical CWE
classification, holding the policy LLM and prompt fixed and comparing seven
strategies at equal token budgets. By isolating policy strength and reward
fidelity, we separate the value of \emph{search structure} from the value of
\emph{extra compute} --- a confound that pervades reported MCTS-for-LLM gains.

Searching a free-form reasoning trajectory beats searching the output taxonomy by
$+12$ leaf points at \emph{half} the token budget on the weak policy
($p{<}10^{-4}$), and on four of five policies overall, because taxonomy search
commits to coarse categories early and compounds the error---making the
\emph{search space} the primary design lever, above budget or reward. The margin
closes, and at the premium end reverses, as the policy strengthens: a policy that
reasons well unaided no longer needs the scaffold. Indeed, the reasoning
strategy's accuracy is flat in its iteration knob: the scaffold (reason, then
commit to a leaf directly), not the tree search on top, carries the gain. At matched compute, reward-guided tree search also delivers a clear gain on
a weak policy ($+8.2$ leaf points over greedy, $+6.4$ over self-consistency) but a
negligible one on a strong policy ($+1.9$); self-consistency---sampling without a
reward---helps neither. The marginal value of search thus \emph{decreases} with
policy strength, contradicting the intuition that better policies search better.
The binding constraint is the accuracy ceiling of the policy on the task, not the
search budget: once a policy nears that ceiling, additional search reconfirms or
drifts. We read this as a budget-aware rule (\S\ref{sec:analysis}): spend on search
when the policy is weak and the reward is informative; spend on a better policy
otherwise.

\paragraph{A budget-aware decision rule.}
Our main practical deliverable is a rule for \emph{when} to spend on search:
read off the policy's pillar accuracy $\polacc$ and the available budget
$\budget$, then choose greedy, self-consistency, or MCTS accordingly
(\S\ref{sec:analysis}).

\paragraph{Future work.}
A trained process reward model (PRM) could shift the reward-bottleneck
conclusion; we left it out to keep the study GPU-free, and flag it as the most
informative next experiment. Extending the protocol to multi-language data and to
detection/localization (beyond classification) are natural follow-ups.

\paragraph{Broader Impact.}
LLM-assisted vulnerability triage is already deployed under tight cost budgets;
our results give practitioners a calibrated answer to ``is structured search
worth its tokens here?''---often the answer is to fix the input representation
and the search space before buying more compute, which lowers both cost and the
energy footprint of security tooling. At the same time, automated CWE labels are
a triage aid, not a verdict: our best configuration still mislabels most
samples at the leaf level, and a false negative can hide a real weakness.
Deployments should keep a human reviewer in the loop and treat predicted labels
as prioritization signals rather than ground truth. The study uses only public,
already-disclosed vulnerabilities and fixes, so it does not create new
disclosure risk.
